\chapter{ЗОРИЛГО}

Энэхүү бакалаврын судалгааны ажлын зорилго нь их сургуулийн SISI портал дээрх оюутанд хамаарах өгөгдөл, үйлчилгээг байгалийн хэлний интерфейсээр дамжуулан хүргэж, оюутны өдөр тутмын ажлыг хөнгөвчлөх хиймэл оюуны туслах системийн загварчлал, архитектур, эхний шатны хэрэгжилтийг боловсруулахад оршино.

Уг зорилгын хүрээнд систем нь оюутанд SISI портал дахь олон цэс, тархай мэдээллийг нэг цэгээс асууж авах боломж олгох, хичээлийн хуваарь, мэдэгдэл, дүн, санхүүгийн мэдээлэл зэрэг давтамж өндөртэй өгөгдөлд хурдан хүрэх замыг бүрдүүлэх, SISI болон цаашид холбогдох бусад MCP үйлчилгээний хооронд давтагддаг гар ажиллагааг бууруулах, мөн оюутны болон багшийн хүсэлтийг системийн бодит үйлдэлд холбох агентын суурь бий болгохыг зорьж байна.

\section{Зорилт}

Дээрх зорилгыг хэрэгжүүлэхийн тулд дараах зорилтуудыг дэвшүүлэв. Эхлээд оюутны портал ашиглах үед тулгардаг давтагддаг асуудал, хэрэгцээг тодорхойлох шаардлагатай. Үүний зэрэгцээ хиймэл оюуны агент, MCP, tool-calling загварын онолын үндсийг судалж, ижил төстэй боловсролын болон байгууллагын туслах системүүдийг харьцуулан үзэх нь чухал. Мөн SISi системийн ухаалаг агентын гол хэрэглэгчид, үндсэн ашиглалтын хувилбарууд, шаардлагуудыг тодорхойлох, системийн өндөр түвшний архитектур болон модулийн зохиомжийг боловсруулах, одоогийн кодын сан дээр тулгуурлан Chrome extension, AI сервер, локал MCP серверийн эхний хэрэгжилтийг баримтжуулах, access token-д суурилсан session reuse болон Chrome extension-ийн ашиглалтын загварыг тодорхойлох, цаашид хөгжүүлэх ерөнхий MCP өргөтгөлүүд, submit/export төрлийн функцүүд, мөн багшийн порталын Excel-д суурилсан дүн оруулах урсгалын төлөвлөгөөг гаргах зорилтуудыг дэвшүүлэв.
